\chapter*{\RL{ملخص}}
\addcontentsline{toc}{chapter}{Arabic Abstract}

\begin{RLtext}
\noindent
في المشهد الاقتصادي الحالي، تفضل الشركات الحديثة اعتماد نظام \L{ERP} موحد قادر على إدارة وتنسيق كافة وظائفها التشغيلية، ولا سيما إدارة الرأسمال البشري (\L{HCM}). وفي الوقت ذاته، تفتح التطورات المتسارعة للذكاء الاصطناعي (\L{AI}) وقدرته على تحليل آلاف الوثائق المعقدة آفاقا جديدة لتطوير عمليات التوظيف وانتقاء أفضل الكفاءات.

\vspace{0.3cm}
يقدم هذا التقرير، المنجز بشركة \L{Cognitus Corporation}، حصيلة تدريب المهندس ضمن مسلك \L{ASEDS} بالمعهد الوطني للبريد والمواصلات (\L{INPT}). وباعتباره أول تجربة لي في أنظمة المؤسسات، بدأت هذه المرحلة بتكوين مكثف في تقنيات \L{SAP} الحديثة من خلال دورات رسمية في منصة \L{BTP} وبرمجة \L{ABAP} ونموذج \L{RAP}، مما مكنني من تطوير مشروع تطبيقي متكامل وفق معايير \L{SAP} الحالية.

\vspace{0.3cm}
يربط النظام المطور بين التوظيف الخارجي وإدارة المؤسسة عبر بنية مزدوجة: بوابة \L{Web} حديثة للمترشحين، متصلة بنظام \L{SAP} الداخلي. وتتيح هذه المنظومة للمسؤولين نمذجة الهيكل التنظيمي، ونشر عروض العمل، وتتبع المترشحين عبر مراحل تقييم محددة. وبفضل محرك الذكاء الاصطناعي لتحليل السير الذاتية وتطويره باستخدام \L{SAP ABAP RAP}، يضمن النظام مسارا سلسا ومؤتمتا من الانتقاء الأولي حتى تسجيل الموظف في \L{SAP}.
\end{RLtext}

\noindent\rule[2pt]{\textwidth}{0.5pt}

\begin{RLtext} 
{\textbf{الكلمات المفتاحية :}}
نظام \L{ERP}، إدارة الرأسمال البشري (\L{HCM})، \L{SAP BTP}، \L{ABAP RAP}، الذكاء الاصطناعي، التوظيف الذكي، الهيكل التنظيمي، \L{React}.
\\
\end{RLtext}

\noindent\rule[2pt]{\textwidth}{0.5pt}